\section{Historia de la evolución técnica: de Logical Agent a Language Agent}

Su Yu (苏煜) divide la historia técnica de los Agent en varias etapas: los primeros logical agent, los neural agent posteriores al año 2000, la rama de semantic parsing, que va del lenguaje a la semántica formal, y los language agent de los últimos tres años. Esta genealogía importa porque evita ver el auge de los Agent en 2026 como algo nuevo surgido de la nada.

\begin{figure}[H]
\centering
\includegraphics[width=0.96\textwidth]{agent-history-timeline.png}
\caption{Genealogía técnica de los Agent: de la lógica simbólica al OpenClaw Moment.}
\end{figure}

\begin{knowledgebox}{Lectura de la figura: la historia técnica no es una historia lineal de victorias}
La línea de tiempo comienza con los logical agent, pasa por los neural agent y el semantic parsing, y termina en los language agent. La tendencia central no es que «los métodos nuevos sustituyen por completo a los antiguos», sino que los LLM convierten el lenguaje en una nueva interfaz general, de modo que los problemas que antes estaban dispersos entre la lógica, la planeación, la invocación de herramientas y la interacción con el entorno vuelven a converger.
\end{knowledgebox}

\subsection{Logical Agent: lógica simbólica y sistemas expertos}

Los Logical Agent eran la visión dominante de la IA temprana: describir el mundo con lógica formal, reglas y sistemas expertos, y luego tomar decisiones mediante un sistema de inferencia. Su ventaja es que son interpretables y tienen una estructura clara; su debilidad es que el mundo real es demasiado complejo, y la cobertura de reglas, el manejo de excepciones y la entrada perceptual son difíciles de escalar. Su Yu enfatiza que perseguir un Agent de inteligencia artificial completo resultaba algo excesivo con las condiciones técnicas de la época, y que eso terminó por dividir la IA en subcampos como la visión, el procesamiento de lenguaje natural y el razonamiento lógico.

\noindent\begin{tabularx}{\textwidth}{p{0.22\textwidth}p{0.34\textwidth}X}
\toprule
Etapa & Mecanismo & Limitaciones \\
\midrule
Logical Agent & Lenguajes lógicos, sistemas expertos, planeación simbólica & Interpretables pero frágiles; dependen de reglas escritas a mano y difícilmente cubren entornos abiertos. \\
Neural Agent & Redes neuronales, aprendizaje por refuerzo, modelos de percepción & Mayor capacidad perceptual, pero carecen de una interfaz de lenguaje general y de capacidad para usar herramientas complejas. \\
Semantic Parsing & Del lenguaje natural a SQL, expresiones lógicas, llamadas a API & Eficaz en entornos específicos, pero cada entorno debe modelarse por separado. \\
Language Agent & LLM como modelo del mundo lingüístico e interfaz de herramientas & Gran capacidad de generalización, pero reliability, memory, costo y aprendizaje continuo siguen siendo cuellos de botella. \\
\bottomrule
\end{tabularx}

\subsection{Semantic Parsing: la prehistoria de los Language Agent}

El Semantic Parsing intenta mapear el lenguaje natural a acciones formalizadas, por ejemplo consultas SQL, consultas a grafos de conocimiento o llamadas a API. Se parece mucho a los Language Agent modernos: ambos deben convertir el lenguaje en conductas ejecutables. La diferencia es que, antes de los LLM, los sistemas por lo general solo funcionaban en una base de datos, un sitio web o un grafo de conocimiento específicos; con los LLM, el modelo incorpora conocimientos previos lingüísticos y conocimiento del mundo más sólidos, lo que le permite generar conductas razonables en muchos más entornos.

\begin{knowledgebox}{Términos clave: Semantic Parsing y Language Agent}
El Semantic Parsing resuelve el problema de «convertir una oración en una forma ejecutable», por ejemplo, transformar «buscar la población de cierta ciudad» en SQL. Los Language Agent extienden esta idea a entornos abiertos: el modelo no solo genera semántica formal, sino que también planea, invoca herramientas, observa resultados y sigue actuando. No hay una ruptura entre ambos, sino una relación de prehistoria y extensión.
\end{knowledgebox}

\subsection{Language Agent: por qué el lenguaje es un acelerador}

Su Yu usa la evolución humana como analogía del papel del lenguaje: el lenguaje apareció muy tarde en la historia evolutiva del ser humano, pero aceleró enormemente el desarrollo de la civilización. De manera similar, los LLM dan a los Agentes de IA una interfaz simbólica general. El lenguaje natural, los lenguajes de programación, los diagramas y los gestos pueden considerarse formas amplias de language; un programming language es solo un tipo de language más formalizado y más adecuado para que una máquina lo ejecute.

\begin{importantbox}{Tesis central: el lenguaje no es una interacción superficial, sino una interfaz con el modelo del mundo}
Con los LLM, los Agent ya no necesitan construir desde cero un analizador semántico para cada entorno. El lenguaje se convierte en la capa intermedia que conecta objetivos, planes, herramientas, retroalimentación y memoria. Coding es importante precisamente porque es uno de los lenguajes formales más potentes del digital world.
\end{importantbox}

\subsection{Resumen de la sección}

La historia técnica de los Agent no comenzó en 2026. Lo nuevo es que los LLM convierten el language en una interfaz general, de modo que los objetivos de los logical agent, la ejecutabilidad del semantic parsing, la capacidad perceptual de los neural agent y la invocación de herramientas vuelven a converger.
